\documentclass[border=6pt]{standalone}
\usepackage{fontspec}
\usepackage{xeCJK}
\usepackage{xcolor}
\usepackage{tikz}
\usetikzlibrary{
    arrows.meta,
    backgrounds,
    calc,
    decorations.pathreplacing,
    fit,
    positioning,
    shapes.geometric
}

\setmainfont{Times New Roman}
\setsansfont{Noto Sans SC}
\setCJKmainfont{Noto Sans SC}
\setCJKsansfont{Noto Sans SC}

\definecolor{ink}{HTML}{1F2937}
\definecolor{muted}{HTML}{64748B}
\definecolor{lineblue}{HTML}{2563EB}
\definecolor{prefill}{HTML}{E7F0FF}
\definecolor{prefillLine}{HTML}{2F5597}
\definecolor{decode}{HTML}{EAF7EA}
\definecolor{decodeLine}{HTML}{2F7D32}
\definecolor{kv}{HTML}{FFF3D8}
\definecolor{kvLine}{HTML}{B7791F}
\definecolor{accent}{HTML}{FDE8E8}
\definecolor{accentLine}{HTML}{C2410C}
\definecolor{panel}{HTML}{F8FAFC}
\definecolor{panelLine}{HTML}{94A3B8}
\definecolor{violet}{HTML}{F1ECFF}
\definecolor{violetLine}{HTML}{6D28D9}
\definecolor{teal}{HTML}{DFF7F3}
\definecolor{tealLine}{HTML}{0F766E}

\tikzset{
    every picture/.style={font=\sffamily},
    title/.style={font=\sffamily\bfseries\Large, text=ink},
    subtitle/.style={font=\sffamily\small, text=muted, align=center},
    block/.style={
        draw=panelLine,
        rounded corners=5pt,
        fill=white,
        line width=0.8pt,
        align=center,
        inner xsep=7pt,
        inner ysep=6pt
    },
    pblock/.style={block, fill=prefill, draw=prefillLine},
    dblock/.style={block, fill=decode, draw=decodeLine},
    kvblock/.style={block, fill=kv, draw=kvLine},
    warnblock/.style={block, fill=accent, draw=accentLine},
    vblock/.style={block, fill=violet, draw=violetLine},
    tblock/.style={block, fill=teal, draw=tealLine},
    cluster/.style={
        draw=panelLine,
        rounded corners=6pt,
        fill=panel,
        line width=0.8pt,
        inner sep=8pt
    },
    arrow/.style={-{Latex[length=3mm,width=2mm]}, draw=ink, line width=0.85pt},
    bluearrow/.style={arrow, draw=prefillLine},
    greenarrow/.style={arrow, draw=decodeLine},
    kvarrow/.style={arrow, draw=kvLine},
    orangearrow/.style={arrow, draw=accentLine},
    control/.style={-{Latex[length=2.6mm,width=1.8mm]}, draw=muted, line width=0.75pt, dashed},
    label/.style={font=\sffamily\scriptsize, text=muted, align=center},
    small/.style={font=\sffamily\scriptsize, align=center},
    rank/.style={
        draw=panelLine,
        rounded corners=3pt,
        fill=white,
        minimum width=12mm,
        minimum height=6mm,
        font=\sffamily\scriptsize,
        align=center
    },
    head/.style={
        draw=ink,
        line width=0.35pt,
        minimum width=6.5mm,
        minimum height=4.8mm,
        font=\sffamily\tiny,
        align=center
    },
    metric/.style={
        draw=panelLine,
        rounded corners=4pt,
        fill=white,
        font=\sffamily\scriptsize,
        align=center,
        inner xsep=5pt,
        inner ysep=4pt
    }
}


\begin{document}
\begin{tikzpicture}[x=1cm,y=1cm]
    \node[title] at (7.9,8.15) {Decode 节点动态 CP 与跨 DP 调度};
    \node[subtitle] at (7.9,7.72) {短请求保持 CP=1 的低开销路径，超长请求按需跨多个 DP-attention 副本协同处理，缓解 Decode 阶段 KV 访存 straggler};

    \node[block, minimum width=2.8cm, minimum height=5.2cm] (traffic) at (1.85,4.25) {};
    \node[font=\sffamily\bfseries\small] at (1.85,6.35) {在线请求流};
    \foreach \i/\y in {0/5.65,1/5.15,2/4.65,3/4.15} {
        \node[metric, fill=white, minimum width=1.45cm] at (1.35,\y) {短请求};
        \node[small, text=muted] at (2.52,\y) {4K--32K};
    }
    \node[metric, fill=accent, draw=accentLine, minimum width=1.45cm] at (1.35,3.25) {长请求};
    \node[small, text=accentLine] at (2.55,3.25) {256K--1M};
    \node[small, text=muted] at (1.85,2.28) {混部流量导致\\DP 副本负载倾斜};

    \node[vblock, minimum width=3.45cm, minimum height=2.55cm] (scheduler) at (5.25,4.9) {};
    \node[font=\sffamily\bfseries\small, text=violetLine] at (5.25,5.9) {CrossDPScheduler};
    \node[small] at (5.25,5.42) {上下文长度};
    \node[small] at (5.25,5.0) {KV Pool 压力};
    \node[small] at (5.25,4.58) {并发与队列深度};
    \node[small, text=violetLine] at (5.25,4.03) {选择 CP 粒度};

    \node[dblock, minimum width=3.55cm, minimum height=4.8cm] (short) at (9.0,4.95) {};
    \node[font=\sffamily\bfseries\small, text=decodeLine] at (9.0,6.72) {短请求路径};
    \node[metric, fill=decode, minimum width=2.0cm] at (9.0,6.1) {CP=1};
    \node[rank, fill=white, minimum width=1.3cm] (dp0) at (8.25,5.3) {DP0};
    \node[rank, fill=white, minimum width=1.3cm] (dp1) at (9.75,5.3) {DP1};
    \node[head, fill=kv, minimum width=1.15cm] at (8.25,4.55) {本地 KV};
    \node[head, fill=kv, minimum width=1.15cm] at (9.75,4.55) {本地 KV};
    \node[small, text=decodeLine] at (9.0,3.55) {低通信开销\\保持常规 DP-attention};

    \node[tblock, minimum width=4.45cm, minimum height=4.8cm] (long) at (13.65,4.95) {};
    \node[font=\sffamily\bfseries\small, text=tealLine] at (13.65,6.72) {长请求路径};
    \node[metric, fill=teal, draw=tealLine, minimum width=2.0cm] at (13.65,6.1) {CP>1};
    \foreach \i/\x in {2/12.3,3/13.2,4/14.1,5/15.0} {
        \node[rank, fill=white, minimum width=0.72cm] (dp\i) at (\x,5.3) {DP\i};
        \node[head, fill=kv, minimum width=0.68cm] at (\x,4.55) {KV};
    }
    \node[metric, fill=white, minimum width=3.45cm] (mgr) at (13.65,3.78) {CrossDPKVCacheManager};
    \node[metric, fill=white, minimum width=3.45cm] (agg) at (13.65,2.95) {score/LSE 聚合};

    \node[block, minimum width=2.3cm, minimum height=1.05cm] (out) at (17.35,4.95) {稳定 Decode\\token 输出};

    \draw[arrow] (traffic.east) -- node[label, above] {请求元数据} (scheduler.west);
    \draw[greenarrow] (scheduler.east) -- node[label, above] {短请求} (short.west);
    \draw[orangearrow] (scheduler.east) to[out=-15,in=180] node[label, below] {长请求} (long.west);
    \draw[greenarrow] (short.east) -- (out.west);
    \draw[greenarrow] (long.east) -- (out.west);

    \node[warnblock, minimum width=4.35cm, minimum height=0.95cm] at (5.25,1.45) {静态 CP：短请求多付通信开销\\或长请求集中成 straggler};
    \node[tblock, minimum width=6.0cm, minimum height=0.95cm] at (12.15,1.45) {动态 CP：用跨 DP 通信换取 HBM 访存分摊\\1M、16 rps 前期实验 P90 ITL 约 60 ms};
\end{tikzpicture}
\end{document}
